\vspace{-1cm}
\section{Giới thiệu đề tài}
\indent Trong bối cảnh nông nghiệp công nghệ cao ngày càng phát triển, đặc biệt là mô hình nhà kính giúp con người kiểm soát chặt chẽ các yếu tố như ánh sáng, nhiệt độ, độ ẩm và chất lượng không khí, việc tối ưu hóa quy trình tưới tiêu vẫn luôn là thách thức. Chúng tôi đề xuất phát triển một hệ thống giám sát và chăm sóc cây trồng thông minh tích hợp trí tuệ nhân tạo (YoloFarm) nhằm xây dựng một giải pháp nông nghiệp thông minh toàn diện, không chỉ dừng lại ở việc giám sát thông số môi trường mà còn có khả năng thấu hiểu sức khỏe cây trồng thông qua AI để đưa ra các biện pháp chăm sóc chủ động.\\
\indent Về công nghệ, chúng tôi sử dụng mạch YOLO:bit làm trung tâm thu thập dữ liệu và điều khiển các thiết bị ngoại vi. Đồng thời, tích hợp trí tuệ nhân tạo để dự báo các biến động môi trường dựa trên dữ liệu lịch sử, nhận diện sâu bệnh dựa trên hình ảnh để đưa ra các cảnh báo sớm.